% \documentclass{article}
% \usepackage[utf8]{inputenc}
% \usepackage{amsmath}
% \usepackage{graphicx}
% \usepackage{booktabs}
% \usepackage{array}
% \usepackage{geometry}\geometry{a4paper, margin=1in}
% \usepackage{hyperref}
% \usepackage{parskip}

% \title{Architectural Description: Human Benchmark \\ \Large Concurrency View}
% \author{Generated by AI Software Architect Assistant}
% \date{\today}

% \begin{document}

% \maketitle
% \tableofcontents

\section{Introduction to the Concurrency View}

The Concurrency View for the Human Benchmark system describes how the system 
manages concurrent execution through its runtime elements. 
It captures the process and thread structure of the system, coordination mechanisms, 
and strategies for ensuring data integrity during concurrent operations. 
The Concurrency View is particularly important for Human Benchmark as it is a web-based 
application supporting real-time multiplayer interactions, simultaneous game sessions, 
and a reactive chat system.

According to the architectural principles discussed in the Software System Architecture Book, 
the Concurrency View helps stakeholders understand how the system handles parallel activities, 
manages shared resources, coordinates activities between runtime components, and scales under load. 
It addresses one of the fundamental aspects of system behavior: how components interact with each 
other at runtime when multiple execution threads exist.

The primary stakeholders for this view are:
\begin{itemize}
    \item \textbf{Developers:} Who need to understand process/thread structure, synchronization mechanisms, and how to avoid race conditions, deadlocks, and resource contention
    \item \textbf{Testers:} Who need to create tests that verify concurrent behavior and simulate simultaneous user interactions
    \item \textbf{Administrators:} Who need to understand startup/shutdown sequences, task failure implications, and system monitoring requirements
    \item \textbf{Communicators:} Who need to explain the system's concurrent behavior to other stakeholders
\end{itemize}

\section{Core Architectural Elements of the Concurrency View}

\subsection{Process and Thread Structure}

The Human Benchmark system employs a multi-tier architecture with several distinct processes
 working together to provide the system's functionality. The concurrency structure is designed to 
 support both single-player and multiplayer.

\subsubsection{Major Process Groups}

\begin{itemize}
    \item \textbf{Client Process Group:} Browser-based React applications running on user devices
    \item \textbf{API Server Process Group:} ASP.NET Core processes handling HTTP requests and WebSocket connections
    \item \textbf{Database Process Group:} PostgreSQL database processes managing data persistence
\end{itemize}

Each process group consists of one or more OS-level processes that work together to fulfill the responsibilities of that group.

% % TODO: DIAGRAM - High-level process group organization with interprocess communication paths
% % SUGGESTED_TYPE: UML Component Diagram with stereotypes for process groups
% % KEY_ELEMENTS_TO_DEPICT: Client Process Group, API Server Process Group, Database Process Group, HTTP/REST communication paths, WebSocket communication paths, Database connection paths
% \begin{figure}[h!]
% \centering
% % \includegraphics[width=0.9\textwidth]{images/placeholder_diagram_for_concurrency_view_process_groups.png}
% \caption{Illustrative Diagram: High-Level Process Groups and Communication Paths (Conceptual Placeholder)}
% \label{fig:concurrency_process_groups}
% \textit{Comment: This diagram should illustrate the three major process groups (Client, API Server, Database) and their communication paths.}
% \end{figure}

\subsubsection{API Server Process Details}

The API Server process group manages several thread pools and execution contexts:

\begin{itemize}
    \item \textbf{HTTP Request Handler Threads:} Process incoming HTTP requests for game configurations, user management, and score recording
    \item \textbf{WebSocket Connection Manager:} Maintains persistent WebSocket connections for real-time communication
    \item \textbf{Game Room Manager Threads:} Handle the creation, management, and monitoring of active game rooms
    \item \textbf{Chat System Threads:} Process and distribute chat messages across different scopes (global, room, group, private)
    \item \textbf{Database Connection Pool:} Manages database connections for data persistence operations
\end{itemize}

% % TODO: DIAGRAM - Detailed thread structure within the API Server Process Group
% % SUGGESTED_TYPE: UML Component Diagram with stereotypes for threads and thread pools
% % KEY_ELEMENTS_TO_DEPICT: HTTP Request Handler Threads, WebSocket Connection Manager, Game Room Manager Threads, Chat System Threads, Database Connection Pool, IPC mechanisms between threads
% \begin{figure}[h!]
% \centering
% % \includegraphics[width=0.9\textwidth]{images/placeholder_diagram_for_concurrency_view_api_server_threads.png}
% \caption{Illustrative Diagram: Thread Structure within API Server Process Group (Conceptual Placeholder)}
% \label{fig:concurrency_api_server_threads}
% \textit{Comment: This diagram should illustrate the different thread pools and execution contexts within the API Server Process Group and their interactions.}
% \end{figure}

% \subsection{Mapping Functional Elements to Concurrency Units}

% This section describes how the functional elements of Human Benchmark map to 
% the system's concurrency units (processes and threads).

% \subsubsection{User Management Functional Elements}

% User management functionality (registration, authentication, profile management) maps to:
% \begin{itemize}
%     \item HTTP Request Handler threads in the API Server Process Group
%     \item Database Connection Pool threads for user data persistence
% \end{itemize}

% These operations are stateless and do not require specialized concurrency considerations beyond standard connection pooling and thread safety practices.

% \subsubsection{Game Execution Functional Elements}

% Game execution functionality maps to:
% \begin{itemize}
%     \item \textbf{Single-player games:} HTTP Request Handler threads for game initialization and score submission
%     \item \textbf{Multiplayer games:} Game Room Manager threads and WebSocket Connection Manager threads that maintain the game state and synchronize player actions in real-time
% \end{itemize}

% \subsubsection{Real-time Communication Functional Elements}

% The real-time communication functionality (multiplayer coordination, chat) maps to:
% \begin{itemize}
%     \item WebSocket Connection Manager threads for maintaining persistent connections
%     \item Game Room Manager threads for game-specific communication
%     \item Chat System threads for processing and routing messages across different scopes
% \end{itemize}

% \subsubsection{Score and Leaderboard Functional Elements}

% Score tracking and leaderboard functionality maps to:
% \begin{itemize}
%     \item HTTP Request Handler threads for score submission and leaderboard queries
%     \item Database Connection Pool threads for score persistence and leaderboard aggregation
%     \item WebSocket Connection Manager threads for real-time leaderboard updates in multiplayer sessions
% \end{itemize}

\subsection{Interprocess Communication (IPC) Mechanisms}

Human Benchmark employs several IPC mechanisms to enable communication between its concurrency units:

\subsubsection{HTTP/REST Communication}

Used for:
\begin{itemize}
    \item Client-to-server communication for non-real-time operations
    \item User authentication and authorization
    \item Game configuration retrieval
    \item Score submission for single-player games
    \item Leaderboard retrieval
\end{itemize}

\subsubsection{WebSocket Communication}

Used for:
\begin{itemize}
    \item Real-time player status updates
    \item Synchronizing game state across multiple clients
    \item Transmitting player actions in multiplayer games
    \item Distribution of chat messages
    \item Live score updates during multiplayer competitions
\end{itemize}

\subsubsection{Database Connections}

Used for:
\begin{itemize}
    \item Persisting user data, game configurations, and scores
    \item Retrieving leaderboard information and historical data
    \item Managing private game access control
\end{itemize}

\subsubsection{Internal Message Queuing}

Used for:
\begin{itemize}
    \item Processing chat messages asynchronously
    \item Managing game room events (player joining/leaving, game state transitions)
    \item Handling WebSocket connection events
\end{itemize}

% % TODO: DIAGRAM - IPC Mechanisms between different components of the system
% % SUGGESTED_TYPE: UML Component Diagram or custom communication diagram
% % KEY_ELEMENTS_TO_DEPICT: Client components, API Server components, Database, HTTP/REST connections, WebSocket connections, Database connections, Internal message queuing
% \begin{figure}[h!]
% \centering
% % \includegraphics[width=0.9\textwidth]{images/placeholder_diagram_for_concurrency_view_ipc_mechanisms.png}
% \caption{Illustrative Diagram: IPC Mechanisms in the Human Benchmark System (Conceptual Placeholder)}
% \label{fig:concurrency_ipc_mechanisms}
% \textit{Comment: This diagram should illustrate the various interprocess communication mechanisms used between components in the Human Benchmark system.}
% \end{figure}

\subsection{State Management}

This section describes the key runtime states of the major functional elements in the Human Benchmark system.

\subsubsection{Game Room State Model}

Game rooms in the multiplayer system transition through several states:

\begin{table}[H]
    \centering
    \begin{tabular}{|c|c|l|}
        \hline
        \textbf{\#} & \textbf{State} & \textbf{Description} \\
        \hline
        1 & \textbf{Created} & Initial state when a room is first created \\
        \hline
        2 & \textbf{Waiting} & Room is open and waiting for players to join \\
        \hline
        3 & \textbf{Ready} & All required players have joined; ready to start \\
        \hline
        4 & \textbf{In Progress} & Game is actively being played \\
        \hline
        5 & \textbf{Round Complete} & Current round of gameplay has finished \\
        \hline
        6 & \textbf{Game Complete} & All rounds have been completed \\
        \hline
        7 & \textbf{Abandoned} & Room has been abandoned or timed out \\
        \hline
    \end{tabular}
    \caption{Game Room State Model}
    \label{tab:game-room-states}
\end{table}


% % TODO: DIAGRAM - State transition diagram for game rooms
% % SUGGESTED_TYPE: UML Statechart Diagram
% % KEY_ELEMENTS_TO_DEPICT: Created, Waiting, Ready, In Progress, Round Complete, Game Complete, and Abandoned states; transitions between states with relevant events/triggers
% \begin{figure}[h!]
% \centering
% % \includegraphics[width=0.8\textwidth]{images/placeholder_diagram_for_concurrency_view_game_room_states.png}
% \caption{Illustrative Diagram: Game Room State Transitions (Conceptual Placeholder)}
% \label{fig:concurrency_game_room_states}
% \textit{Comment: This diagram should illustrate the state transitions for multiplayer game rooms in the Human Benchmark system.}
% \end{figure}

\subsubsection{Player Connection State Model}

Each player connection in the system transitions through these states:

\begin{table}[H]
    \centering
    \begin{tabular}{|c|c|l|}
        \hline
        \textbf{\#} & \textbf{State} & \textbf{Description} \\
        \hline
        1 & \textbf{Connected} & Initial state when a WebSocket connection is established \\
        \hline
        2 & \textbf{Authenticated} & User has been authenticated on the connection \\
        \hline
        3 & \textbf{Idle} & Connection is active but not currently in a game \\
        \hline
        4 & \textbf{In Game} & Connection is associated with an active game \\
        \hline
        5 & \textbf{Disconnecting} & Connection is in the process of closing \\
        \hline
        6 & \textbf{Disconnected} & Connection has been terminated \\
        \hline
    \end{tabular}
    \caption{Player Connection State Model}
    \label{tab:player-connection-states}
\end{table}


\subsubsection{Chat Message State Model}

Chat messages in the system transition through these states:

\begin{table}[H]
    \centering
    \begin{tabular}{|c|c|l|}
        \hline
        \textbf{\#} & \textbf{State} & \textbf{Description} \\
        \hline
        1 & \textbf{Created} & Message has been created by sender \\
        \hline
        2 & \textbf{Processing} & System is processing the message (e.g., content filtering, routing) \\
        \hline
        3 & \textbf{Delivered} & Message has been delivered to recipients \\
        \hline
        4 & \textbf{Read} & Message has been read by recipients (for private chats) \\
        \hline
        5 & \textbf{Archived} & Message has been archived in message history \\
        \hline
    \end{tabular}
    \caption{Chat Message State Model}
    \label{tab:chat-message-states}
\end{table}


\subsection{Synchronization and Integrity}

The Human Benchmark system employs several mechanisms to ensure data integrity during concurrent access:

\subsubsection{Game State Synchronization}

\begin{itemize}
    \item \textbf{Centralized Authority:} Server acts as the authoritative source for game state
    % \item \textbf{Optimistic Concurrency:} Client-side prediction with server-side validation and correction
    \item \textbf{Event Sequencing:} Events are timestamped and sequenced to ensure consistent processing order
    \item \textbf{Conflict Resolution:} Server resolves conflicting simultaneous actions based on predefined rules
\end{itemize}

\subsubsection{Database Access Synchronization}

\begin{itemize}
    \item \textbf{Connection Pooling:} Managed by Entity Framework Core to handle concurrent database access
    \item \textbf{Transactions:} Used for multi-step database operations that must be atomic
    \item \textbf{Row-level Locking:} Implemented for critical data updates to prevent conflicts
    \item \textbf{Optimistic Concurrency Control:} For less critical data, using row version/timestamp checking
\end{itemize}

\subsubsection{Memory Access Synchronization}

\begin{itemize}
    \item \textbf{Thread-safe Collections:} Used for shared data structures accessed by multiple threads
    \item \textbf{Immutable Objects:} Where appropriate, to eliminate the need for locks
    \item \textbf{Lock Mechanisms:} Carefully applied where necessary for shared mutable state
    \item \textbf{Asynchronous Programming:} Used extensively with async/await pattern to avoid blocking threads
\end{itemize}

% idk, feelt pointless
% \subsection{Startup and Shutdown Procedures}

% \subsubsection{System Startup Sequence}

% The startup sequence for the Human Benchmark concurrency units follows this order:

% \begin{enumerate}
%     \item Database services initialize first
%     \item API Server services start and establish database connections
%     \item WebSocket services initialize and prepare for client connections
%     \item Internal message processing systems start
%     \item Game Room Manager initializes (but creates no rooms yet)
%     \item System signals readiness for client connections
% \end{enumerate}

% \subsubsection{System Shutdown Sequence}

% The shutdown sequence follows this order:

% \begin{enumerate}
%     \item New client connections are rejected
%     \item Active game rooms are notified of impending shutdown
%     \item Client connections are gracefully terminated with appropriate messages
%     \item Game Room Manager shuts down all active rooms
%     \item WebSocket connections are closed
%     \item Internal message queues are drained and processors shut down
%     \item Database connections are closed
%     \item API Server services shut down
%     \item Database services shut down
% \end{enumerate}

\subsection{Fault Tolerance and Recovery}

\subsubsection{Process Failure Handling}

\begin{itemize}
    \item \textbf{Client Process Failures:} Detected through WebSocket heartbeats; player state preserved for reconnection
    \item \textbf{API Server Process Failures:} Handled by Kubernetes auto-restart and session persistence
    \item \textbf{Database Process Failures:} Managed through PostgreSQL's built-in replication and failover
\end{itemize}

\subsubsection{Recovery Mechanisms}

\begin{itemize}
    \item \textbf{Connection Recovery:} Clients can reconnect and resume sessions after brief disconnections
    \item \textbf{State Recovery:} Game state is periodically persisted to allow restoration
    \item \textbf{Partial Failure Mode:} System can operate with reduced functionality if some components fail
\end{itemize}

\subsection{Scalability Strategy}

The Human Benchmark system's concurrency design supports scalability through:

\begin{itemize}
    \item \textbf{Stateless API Design:} HTTP endpoints are designed to be stateless for horizontal scaling
    \item \textbf{WebSocket Connection Sharding:} WebSocket connections can be distributed across multiple servers
    \item \textbf{Database Connection Pooling:} Efficient use of database connections
    \item \textbf{Container Orchestration:} Kubernetes enables dynamic scaling of API Server processes
    \item \textbf{Game Room Partitioning:} Game rooms can be distributed across server instances
\end{itemize}

\section{Perspective Priorities for the Concurrency View}

This section analyses the relevance of various architectural perspectives to the Concurrency View of the Human Benchmark system.

\begin{table}[H]
\centering
\caption{Perspective Priorities for the Concurrency View of Human Benchmark}
\label{tab:perspective_priorities_concurrency}
\begin{tabular}{p{0.3\textwidth}p{0.15\textwidth}p{0.45\textwidth}}
% \toprule
\textbf{Architectural Perspective} & \textbf{Priority} & \textbf{Rationale for Human Benchmark within this View} \\
% \toprule
Accessibility & \greencircle Low & Handled mainly by the UI layer, little impact. \\
Availability and Resilience & \redcircle High & The real-time multiplayer nature of Human Benchmark requires high availability of concurrent processes and resilience to failures to maintain game integrity. \\
Development Resource & \yellowcircle Medium & Threading models and concurrency patterns require careful design and implementation, impacting development effort and requiring specialized skills. \\
Evolution & \yellowcircle Medium & The concurrency design must support future additions of new game types and scaling to more concurrent users. \\
Internationalization & \greencircle Low & Mostly a UI/content concern, not concurrency related. \\
Location & \yellowcircle Medium & WebSocket latency between distributed players can affect gameplay experience, and the system must handle players from different geographical regions. \\
Performance and Scalability & \redcircle High & Real-time interactions in multiplayer games demand high performance under concurrent load, and the system must scale to support many simultaneous game sessions. \\
Regulation & \greencircle Low & While important for the system overall, regulatory concerns have minimal impact on the concurrency mechanisims. \\
Security & \yellowcircle Medium & Concurrent access to shared resources requires careful authentication and authorization to prevent security issues, particularly in multiplayer contexts. \\
Usability & \yellowcircle Medium & Responsiveness during concurrent operations directly affects user experience, particularly in real-time games measuring reaction time. \\
% \toprule
\end{tabular}
\end{table}

\section{Impact of Key Perspectives on the Concurrency View}

Based on the priorities identified above, this section examines how the highest-priority perspectives influence the Concurrency View of the Human Benchmark system.

\subsection{Performance and Scalability Perspective}

Applying the Performance and Scalability perspective to the Concurrency View has significantly shaped several aspects of the architecture:

\begin{itemize}
    \item \textbf{Thread Pool Optimization:} The system uses carefully tuned thread pools in the API Server processes to balance responsiveness with resource utilization. This is particularly important for reaction time tests where milliseconds matter.
    
    \item \textbf{Asynchronous Processing:} Non-critical operations (like score persistence and leaderboard updates) are handled asynchronously to prevent blocking the main game processing threads, ensuring smooth gameplay even under high load.
    
    \item \textbf{Connection Management:} The WebSocket Connection Manager implements connection limiting and prioritization to prevent resource exhaustion during traffic spikes, ensuring that active game sessions maintain quality of service.
    
    \item \textbf{Stateless Design:} Where possible, HTTP endpoints are designed to be stateless, allowing horizontal scaling of API servers without session affinity requirements.
    
    \item \textbf{Database Connection Pooling:} The system carefully manages database connections to prevent connection storms and ensure efficient use of database resources during concurrent gameplay.
\end{itemize}

These design decisions enable Human Benchmark to maintain high performance during peak usage periods and scale to support thousands of concurrent users across multiple game types.

\subsection{Availability and Resilience Perspective}

The Availability and Resilience perspective has influenced the Concurrency View in the following ways:

\begin{itemize}
    \item \textbf{Process Isolation:} Critical system components are isolated in separate processes to contain failures and prevent system-wide outages. For example, the failure of a single game room will not affect other active game rooms.
    
    \item \textbf{State Persistence:} Game state is periodically persisted to enable recovery in case of process failures, allowing users to resume their sessions after brief interruptions.
    
    \item \textbf{Failure Detection:} The system implements heartbeat mechanisms for WebSocket connections to quickly detect client disconnections or API server failures.
    
    \item \textbf{Graceful Degradation:} Under extreme load, the system prioritizes maintaining existing game sessions over accepting new ones, ensuring quality of service for active users.
    
    \item \textbf{Concurrent Recovery:} Recovery operations run concurrently with normal system operation to minimize downtime and service disruption.
\end{itemize}

These resilience strategies ensure that the Human Benchmark system can maintain its critical real-time functionality even in the face of partial component failures or unexpected load spikes.

\subsection{Location Perspective}

The Location perspective has influenced the Concurrency View in the following ways:

\begin{itemize}
    \item \textbf{Connection Latency Management:} The concurrency design accounts for varying network latencies between players in different geographic locations, implementing compensation mechanisms to ensure fair gameplay.
    
    \item \textbf{Optimistic Concurrency Control:} To mitigate the effects of high latency, the system implements client-side prediction with server validation, reducing the perceived impact of network delays.
    
    \item \textbf{Regional Partitioning:} Game rooms can be regionally partitioned to group players with similar latencies together, improving the multiplayer experience.
    
    \item \textbf{Adaptive Timing:} Time-critical game mechanics adapt to measured network conditions to ensure fairness across different connection qualities.
\end{itemize}

These location-aware concurrency strategies enable Human Benchmark to deliver a consistent and fair experience to users regardless of their geographic location.

\section{Importance and Justification of the Concurrency View for Human Benchmark}

The Concurrency View is of critical importance for the Human Benchmark system for several compelling reasons:

\subsection{Real-Time Interaction Requirements}

Human Benchmark is fundamentally designed to measure reaction time and cognitive performance in milliseconds. The system's core value proposition depends on accurate timing and responsive interaction. The Concurrency View is essential for ensuring that the system can:

\begin{itemize}
    \item Process user inputs with minimal latency
    \item Maintain accurate timing for reaction tests
    \item Synchronize game state across multiple clients in real-time
    \item Handle concurrent scoring and leaderboard updates without introducing delays
\end{itemize}

Without a well-designed concurrency architecture, the system would fail to deliver its primary function of providing accurate cognitive and reflex measurements.

\subsection{Multiplayer Competitive Environment}

The multiplayer aspects of Human Benchmark introduce complex concurrency challenges that must be addressed:

\begin{itemize}
    \item Simultaneous actions from multiple players must be fairly processed and resolved
    \item Game state must remain consistent across all participants
    \item Real-time competition requires synchronization of game events and scoring
    \item Elimination mechanics in multiplayer games demand precise timing and state management
\end{itemize}

The Concurrency View directly addresses these challenges by defining how the system handles simultaneous interactions and maintains state consistency across distributed clients.

\subsection{Chat System Requirements}

The recently implemented chat system introduces additional concurrency requirements:

\begin{itemize}
    \item Multiple message scopes (global, room, group, private) require concurrent message routing
    \item Real-time message delivery needs efficient thread management
    \item Message synchronization across multiple recipients must be handled correctly
\end{itemize}

The Concurrency View provides clarity on how these requirements are met through the chat system threads and message processing mechanisms.

\subsection{Scalability and Performance Demands}

As a web-based platform, Human Benchmark must support varying load levels:

\begin{itemize}
    \item Concurrent user sessions ranging from dozens to thousands
    \item Multiple simultaneous multiplayer game rooms
    \item Peak load periods with surges in user activity
\end{itemize}

The Concurrency View details how the system scales to meet these demands through its process structure, thread management, and resource allocation strategies.

% seemed alrady too much
% \subsection{Integration with Architectural Principles}

% The importance of the Concurrency View aligns with key principles from the Software System Architecture Book:

% \begin{itemize}
%     \item "It is not possible to capture all features and properties of a complex system in a single comprehensible model." The Concurrency View captures aspects of the system that would be difficult to represent in other views.
    
%     \item "Architectures are created solely to meet stakeholder needs." For Human Benchmark, stakeholders need a system that can handle concurrent users, provide fair multiplayer competition, and deliver responsive real-time interactions.
    
%     \item "A concern, problem, or system element is architecturally significant if it has a wide impact on the structure of the system or on its important quality properties." The concurrency aspects of Human Benchmark are architecturally significant as they directly impact key quality properties like performance, scalability, and fairness.
% \end{itemize}

% In conclusion, the Concurrency View is not just important but essential for Human Benchmark due to the system's core focus on real-time interaction, competitive multiplayer capabilities, and the need to ensure fairness and accuracy in cognitive measurements across many concurrent users.

\section{Cross-View Consistency Considerations}

The Concurrency View must maintain consistency with other
 architectural views to ensure a coherent overall system architecture. 
 Key cross-view considerations include:

\begin{figure}[H]
    \centering
    \includegraphics[width=0.7\textwidth]{AutomationAndSections/Sections/4concurencyView/Images/ConcurrencyView_crossover.png}
    \caption{Concurrency View Cross-View Consistency diagram}
    \label{fig:concurrency_view_cross_view}
\end{figure}


% \begin{itemize}
%     \item The functional elements identified in the Functional View must be correctly mapped to concurrency units in this view
%     \item The interactions between functional elements must align with the interprocess communication mechanisms described here
%     \item State transitions described in this view must support the functional behaviors described in the Functional View
% \end{itemize}

% \subsection{Consistency with the Deployment View}

% \begin{itemize}
%     \item The process structure described in this view must be deployable on the hardware infrastructure described in the Deployment View
%     \item Scaling strategies for concurrency units must align with the deployment scaling capabilities
%     \item Network communication patterns must be supported by the network infrastructure described in the Deployment View
% \end{itemize}

% \subsection{Consistency with the Information View}

% \begin{itemize}
%     \item Data access patterns and concurrency control mechanisms must support the data models described in the Information View
%     \item State models in this view must be consistent with persistent entity lifecycles in the Information View
%     \item Real-time data flows must align with the information structure and integrity constraints
% \end{itemize}

% \subsection{Consistency with the Development View}

% \begin{itemize}
%     \item Concurrency patterns and mechanisms must be implementable using the technologies and frameworks specified in the Development View
%     \item Module organization must support the concurrency units described in this view
%     \item Development standards must address thread safety requirements identified in this view
% \end{itemize}

% This document forms one part of a comprehensive Architectural Description (AD) for the Human Benchmark system. When implementing changes related to concurrency, architects and developers should consult other relevant views to ensure system-wide consistency and avoid unintended consequences.

% \end{document}